\section{Conclusion}
\label{sec:conclusion}

The present study investigated how the predictive performance of a standalone decision tree compared with decision-tree ensemble methods for ten-class network intrusion detection on a 257,673-row subset of the University of New South Wales network-based 2015 (UNSW-NB15) dataset~\cite{Moustafa2015}. An extended iterative dichotomiser three (ID3) decision-tree learner, bottom-up reduced-error post-pruning, bootstrap aggregation, known as bagging, and stagewise additive modelling using a multi-class exponential loss function (SAMME) were implemented within a unified software framework so that the base classifier family and 41-attribute representation remained constant across all models. Using a fixed, stratified four-way split that separated model growth with 123,683 rows, reduced-error pruning with 41,225 rows, development validation with 41,234 rows, and final generalisation testing with 51,531 rows, four sequential experimental stages were executed to answer the four research questions of the study.

First, regarding sample-size scaling and class scarcity, the Stage~One learning-curve experiments across 2,000 to 50,000 rows showed that increasing the volume of original training data produced larger gains in minority-attack recognition than applying ensemble methods to small subsamples. Under natural class proportions, rare categories such as \emph{Worms}, representing $0.07\%$ of traffic, had too few distinct instances in small samples for bootstrap resampling to overcome. Grouping the training samples by exact 41-attribute vectors further demonstrated that the decrease in unrestricted training accuracy from $99.00\%$ at 2,000 rows to $93.26\%$ at 50,000 rows matched the exact deterministic ceiling imposed by identical predictor vectors carrying conflicting attack labels.

Second, regarding the interaction between structural regularisation and class balancing, Stage~One and Stage~Two showed that unweighted reduced-error pruning behaved differently depending on the training class distribution. When trees were induced on naturally skewed data, post-pruning consistently increased overall accuracy $a$ while reducing macro-averaged F1-score $\bar{f}$ by collapsing sparse minority-class subtrees. Conversely, when minority classes were oversampled during tree induction and the induced tree was post-pruned against a natural-proportion pruning set, class balancing and post-pruning acted complementarily. In the Stage~Three full-pool pilot, hybrid class balancing at $\lambda = 0.5$ outperformed pure equal-quota balancing at $\lambda = 1.0$ across every model family in both validation $\bar{f}$ and overall accuracy $a$.

Third, regarding ensemble behaviour, 25-member bagging consistently improved over unpruned single trees across all tested sample sizes and class policies, raising generalisation $\bar{f}$ from $0.5545$ to $0.5747$ and overall accuracy from $77.44\%$ to $79.55\%$. Extending bagging from 25 to 35 and 50 members on the validation set produced no further gain in $\bar{f}$, measuring $0.6026$ at 25 members versus $0.6004$ at 35 and 50 members. Sequential SAMME boosting benefited substantially from increasing the base-tree depth limit from $d = 3$ to $d = 8$ at learning rate $\eta = 1.0$, but lagged behind both bagging and the unpruned tree on the generalisation set with $\bar{f} = 0.5408$, exhibiting lower precision on \emph{Shellcode} and \emph{Reconnaissance} and lower recall on \emph{Normal} traffic.

Fourth, in the locked full-pool generalisation comparison under the shared $\lambda = 0.5$ hybrid balancing policy, the standalone reduced-error-pruned decision tree achieved the strongest overall performance across all three primary metrics: a mean macro-averaged F1-score of $0.5993$, a balanced accuracy of $63.02\%$, and an overall accuracy of $80.73\%$, while reducing mean leaf count by $93.7\%$ to 1,861 leaves occupying $0.08$ megabytes (MB) and completing batch inference $13.8$ times faster than 25-member bagging. While 25-member bagging achieved the highest F1-score on \emph{Normal}, \emph{Fuzzers}, \emph{Shellcode}, and \emph{Generic}, unpruned bagging generated substantially more false \emph{Backdoor} and \emph{Analysis} predictions on true \emph{DoS} and \emph{Exploits} traffic than the pruned standalone tree.

Four directions for future work arise directly from the findings and limitations of the present study:
\begin{enumerate}
    \item \textbf{Bagging of post-pruned balanced trees:} Because bagging improved F1-score on \emph{Normal}, \emph{Fuzzers}, \emph{Shellcode}, and \emph{Generic} while reduced-error post-pruning suppressed false alarms on oversampled \emph{Backdoor} and \emph{Analysis} leaves, applying reduced-error pruning to each individual balanced bagging member is a natural extension.
    \item \textbf{Macro-averaged F1-score pruning criteria:} Replacing the unweighted error count in reduced-error pruning with a class-weighted error criterion or direct validation $\bar{f}$ optimisation would align the post-pruning objective with balanced multiclass detection.
    \item \textbf{Deduplicated and temporal-split evaluation:} Future experiments on UNSW-NB15 should evaluate deduplicated flow sets and partition records by capture time window or internet protocol (IP) subnet to measure generalisation across distinct network sessions.
    \item \textbf{Random attribute subspace selection:} Extending the bagging implementation to sample a random subset of descriptive attributes at each split node would decorrelate root splits dominated by high-cardinality \texttt{proto} and \texttt{service} attributes.
\end{enumerate}
